\chapter{Conclusion and Future Work}\label{ch:future_work}

\section{Conclusion}

This thesis asked how a cell-free ISAC transmitter can minimize power while
providing robust communication QoS and target-specific sensing service with a
physically implementable AP--target topology. The answer developed here is a
target-centric allocation framework in which communication beamforming remains
globally cooperative, while binary association variables authorize only the
dedicated sensing covariance of each target. This separation aligns the sensing
topology, the target-specific FIM/PCRB requirement, and the dedicated sensing
power that is actually controlled by the optimization.

The resulting mixed-integer formulation jointly accounts for worst-case
communication SINR, sensing SINR, matrix PCRB, dedicated-sensing authorization,
and total AP-power limits. Covariance lifting, the S-Procedure, and a
Schur-complement representation make the robust and tracking constraints
solver-ready. A dual DC-penalty SCA procedure then treats rank-one and binary
deficiencies jointly, while direct validation distinguishes a recovered
physical design from a relaxed optimization point.

The numerical evidence answers the practical part of the question. In the
tested scenarios, topology-aware recovery improves physical feasibility in
ordinary and pressure geometries; robust design trades a modest power increase
for markedly lower observed outage within the CSI-uncertainty model; and the
dual-DC workflow produces small rank and binary deficiencies before physical
validation. The larger-dimensional experiment further shows usable recovery
behaviour at the tested scale, with a computational cost above that of
fixed-topology baselines. These are empirical, configuration-specific findings
rather than guarantees of global optimality or universal relaxation tightness.

\section{Future Work}

The next priority is broader, reproducible validation. Future experiments
should report feasibility, recovered physical transmit power, worst-case
communication SINR, sensing SINR, PCRB, per-AP power, and
association-cardinality residuals for every sweep, while preserving the
distinction between relaxed and recovered solutions. Larger independent seed
sets, confidence intervals, and more systematic PCRB, cluster-size, and
CSI-error sweeps are needed to quantify the robustness--sensing--power
trade-off with greater confidence.

The model can also be extended in directions that change the current
single-interval assumptions: dynamic target-state prediction across scheduling
intervals, uncertainty in sensing channels and target parameters, fronthaul or
processing costs for sensing clusters, and distributed implementations for
larger cell-free networks. Each extension should retain direct validation of
the recovered physical variables so that added modelling flexibility does not
obscure the original communication, sensing, tracking, and AP-power guarantees.
